\label{chap:p3-83-47-godel_cohen}

\section{De la consistencia constructible a la independencia por forcing}

\begin{theorem}[Gödel, 1940: mitad constructible]
En la metateoría usual,
\[
  \operatorname{Con}(\mathsf{ZF})
  \Longrightarrow
  \operatorname{Con}
  \bigl(\mathsf{ZF}+\mathsf{AC}+\mathsf{GCH}\bigr).
\]
Más precisamente, dado un modelo \(M\models\mathsf{ZF}\), su universo
constructible interno \(L^M\) satisface, en el sentido interno correspondiente,
\[
  \mathsf{ZF}+V=L+\mathsf{AC}+\mathsf{GCH}.
\]
\end{theorem}

\begin{proof}[Idea de la prueba clásica]
La jerarquía constructible se define por recursión sobre los ordinales,
admitiendo en cada etapa únicamente los subconjuntos definibles sobre etapas
anteriores. La construcción conserva los ordinales, produce un modelo interno
y proporciona un buen orden definible global, del que se deducen elección y
la hipótesis generalizada del continuo \citep{goedel1940}.
\end{proof}

\begin{theorem}[Cohen, 1963--1964: mitad forzante]
Se tienen las implicaciones metamatemáticas
\[
  \operatorname{Con}(\mathsf{ZFC})
  \Longrightarrow
  \operatorname{Con}(\mathsf{ZFC}+\neg\mathsf{CH}),
\]
y
\[
  \operatorname{Con}(\mathsf{ZF})
  \Longrightarrow
  \operatorname{Con}(\mathsf{ZF}+\neg\mathsf{AC}).
\]
Combinando las dos mitades, \(\mathsf{CH}\) es independiente de
\(\mathsf{ZFC}\) y \(\mathsf{AC}\) es independiente de \(\mathsf{ZF}\),
siempre bajo la hipótesis de consistencia correspondiente
\citep{cohen1963,cohen1964}.
\end{theorem}

Para la primera construcción se parte de un modelo base \(M\), un poset
\(\mathbb P\in M\) y un filtro \(M\)-genérico \(G\subseteq\mathbb P\):
\[
  G\cap D\ne\varnothing
  \qquad
  \text{para todo denso }D\subseteq\mathbb P\text{ perteneciente a }M.
\]
El teorema de forcing relaciona \(p\Vdash_M\varphi\) con la verdad de
\(\varphi\) en \(M[G]\). Para obtener
\(\mathsf{ZF}+\neg\mathsf{AC}\) se toma además un submodelo simétrico
adecuado de una extensión genérica; una extensión ordinaria de un modelo de
\(\mathsf{ZFC}\) conserva elección \citep{cohen1966}.

\begin{remark}[Condiciones de validez]
Estos teoremas prueban consistencia relativa e independencia. No prueban la
consistencia absoluta de \(\mathsf{ZF}\) o \(\mathsf{ZFC}\), no deciden
\(\mathsf{CH}\) en un universo privilegiado y no convierten cualquier filtro
compatible en un filtro genérico. Tampoco seleccionan las ramas HMT ni
identifican una salida numérica.
\end{remark}

\subsection{\(2\) operaciones no equivalentes: extensión del modelo y enriquecimiento del objeto}

El forcing y HMT no realizan la misma operación. El primero compara un
universo \(M\) con una extensión \(M[G]\); la segunda conserva fibras que un
lector olvida. Si \(\mathsf{Hist}_{\rm HMT}\) es la categoría de historias
enriquecidas, el functor
\[
  U:\mathsf{Hist}_{\rm HMT}\longrightarrow\mathsf{Set},
  \qquad
  (x_m,\mu_m)_{m\ge0}\longmapsto(x_m)_{m\ge0},
\]
puede identificar historias distintas sin cambiar el conjunto visible. La
pregunta de Cohen es qué subconjuntos existen en el modelo; la pregunta
genealógica es cuántas elevaciones con memoria y operadores posee un objeto
visible. Son ejes compatibles, pero ninguno sustituye al otro.

\begin{center}
\begin{tabularx}{.94\textwidth}{>{\bfseries}p{.22\textwidth}XX}
\toprule
Operación & Qué cambia & Qué no prueba por sí sola\\
\midrule
\(M\mapsto M[G]\) & Universo de conjuntos, nombres y verdad forzada &
Unicidad o genealogía de una sección HMT.\\
\(\xi\mapsto U(\xi)\) & Cantidad de memoria visible en la proyección &
\(\mathsf{CH}\), \(\mathsf{AC}\) o genericidad sobre \(M\).\\
\bottomrule
\end{tabularx}
\end{center}

\subsection{La potencia HMT no presupone \texorpdfstring{\(\mathsf{CH}\)}{CH}}

En cada modelo donde \(X\) es un espacio de Cantor,
\[
  |\Clop(X)|=\aleph_0,\qquad
  |\Closed(X)|=\mathfrak c,\qquad
  |\mathcal P(X)|=2^{\mathfrak c},
  \qquad \mathfrak c=2^{\aleph_0}.
\]
\(\mathsf{CH}\) afirma únicamente \(\mathfrak c=\aleph_1\). Por el teorema
de Cantor, \(\mathfrak c<2^{\mathfrak c}\) en cualquiera de los modelos
comparados. Al pasar de \(M\) a \(M[G]\) pueden aparecer nuevas ramas, nuevos
códigos cerrados y nuevos subconjuntos, aunque los cilindros finitos
conserven su regla de formación. La potencia plena es sensible al modelo; la
distinción clopen--cerrada--plena no depende de decidir \(\mathsf{CH}\).

\section{Poset de prefijos compatibles y límite inverso de las extensiones}

Definimos
\[
  \mathbb P_{\TPK}=\bigsqcup_{n\ge0}\mathcal S_n
\]
y ordenamos por prolongación: \(q\le p\) si \(q\) contiene más información
que \(p\). Para cada \(N\),
\[
  D_N=\{p\in\mathbb P_{\TPK}:\ell(p)\ge N\}.
\]
Si todo prefijo superviviente puede prolongarse, cada \(D_N\) es denso.

\begin{theorem}[Ramas y filtros de profundidad]
Las secciones de \(\OmegaInf\) están en biyección con los filtros compatibles
de \(\mathbb P_{\TPK}\) que intersectan todos los \(D_N\).
\end{theorem}

\begin{proof}
Los prefijos de una sección son compatibles, dirigidos y alcanzan toda
profundidad. Recíprocamente, un filtro que corta cada \(D_N\) determina una
única restricción en cada nivel; la direccionalidad excluye condiciones
incompatibles y produce una sección.
\end{proof}

Encontrar todos los densos \(D_N\) expresa profundidad ilimitada, no
genericidad sobre un modelo. Ser \(M\)-genérico exigiría intersectar todo
denso \(D\in M\).

\begin{proposition}[Caso Cohen condicionado]
Si \(\OmegaInf\) es compacto, metrizable, cero-dimensional, perfecto y sin
puntos aislados, es homeomorfo al Cantor; el forcing de cilindros no vacíos
tiene una parte densa equivalente al forcing de Cohen
\citep{cohen1963,cohen1964}.
\end{proposition}

Una sección única produce, por el contrario, un forcing atómico o trivial.
La palabra «forcing» no selecciona por sí misma una rama.

Para una sección \(x=(x_m)_{m\ge0}\), sus prefijos forman un filtro \(G_x\)
que atraviesa los densos de profundidad. Sin embargo,
\[
  \bigl(\forall N,\;G_x\cap D_N\ne\varnothing\bigr)
  \centernot\Longrightarrow
  \bigl(\forall D\in M\text{ denso},\;G_x\cap D\ne\varnothing\bigr).
\]
La primera propiedad expresa prolongación ilimitada; la segunda,
genericidad relativa a un modelo. Una sección compatible producida por
nueve fases del TPK no se convierte por nomenclatura en un real de Cohen.

\begin{remark}[Condiciones de validez]
Falsan una identificación automática con Cohen: un cilindro terminal o un
punto aislado; un filtro que atraviesa todos los \(D_N\) pero omite algún
denso de \(M\); dos memorias compatibles sobre la misma rama visible; o la
ausencia de un embebimiento denso explícito. El nombre compartido «rama» no
es un morfismo.
\end{remark}

\section{Correspondencia entre secciones genealógicas y clases de Borel del espacio límite}

Los cilindros generan
\[
 \Gamma_{\rm HMT}^{\rm Bor}
 =
 \sigma\bigl(\{[p]_n:p\in\mathcal S_n\}\bigr).
\]
La clase es estable por complementos, uniones numerables e imágenes inversas
de mapas continuos certificados.

Si un juego infinito sobre historias posee condición de victoria boreliana,
el teorema de determinación boreliana de Martin garantiza una estrategia
ganadora \citep{martin1975}. Ésta es una aplicación clásica al espacio HMT;
no deriva AD, MM, un valor de \(2^{\aleph_0}\) ni categoricidad de
\(\mathcal P(\mathbb R)\).

Para \(A,B\subseteq\OmegaInf\), la reducción de Wadge
\[
 A\le_WB
\]
requiere un mapa continuo \(f\) con \(x\in A\iff f(x)\in B\). Si se
restringen los mapas a operadores TPK certificados, aparece un preorden más
estrecho cuya buena fundación o semilinealidad necesita una prueba propia.

\begin{remark}[Condiciones de validez]
Una demostración que usa la buena fundación Wadge para probar la
determinación que luego justificaría esa buena fundación es circular. Las
ramas históricas AD/MM se conservan como arquitectura, no como teoremas HMT.
\end{remark}

\section{Optimización minimax sobre ventanas finitas de prefijos compatibles}

Sea \(P_N=\mathcal S_N\) el conjunto finito de candidatos, \(T_N\) un conjunto
finito de pruebas adversariales y
\[
  A_N:P_N\times T_N\to[-1,1]
\]
un pago fijado antes del resultado.

\begin{theorem}[Minimax finito]
\[
 \max_{\sigma\in\Prob(P_N)}
 \min_{\tau\in\Prob(T_N)}
 \mathbb E_{\sigma,\tau}A_N
 =
 \min_{\tau\in\Prob(T_N)}
 \max_{\sigma\in\Prob(P_N)}
 \mathbb E_{\sigma,\tau}A_N.
\]
\end{theorem}

\begin{proof}
Es la dualidad fuerte entre el programa lineal del selector y el del
adversario, en la forma del teorema minimax de von Neumann
\citep{vonneumann1928}.
\end{proof}

El equilibrio puede ser mixto y no proporciona una rama pura única.

\section{Sistema proyectivo de truncamientos y construcción de la sección límite}

Sean
\[
 X=\varprojlim P_N,
 \qquad
 Y=\varprojlim T_N
\]
compactos, y sea \(A:X\times Y\to\mathbb R\) continuo. Los espacios de
probabilidades son compactos convexos en la topología débil. El teorema de
Sion \citep{sion1958} da
\[
 \sup_{\mu\in\Prob(X)}\inf_{\nu\in\Prob(Y)}
 \iint A\,d\mu\,d\nu
 =
 \inf_{\nu\in\Prob(Y)}\sup_{\mu\in\Prob(X)}
 \iint A\,d\mu\,d\nu.
\]

\begin{corollary}
Un equilibrio global selecciona una rama pura sólo si la medida del selector
es una Dirac \(\delta_x\). La existencia de equilibrio mixto no implica
concentración.
\end{corollary}

\begin{remark}[Condiciones de validez]
El paso al límite falla sin compacidad, continuidad o consistencia
proyectiva de las medidas. Si el pago se define usando el blanco que debe
seleccionar, el minimax es target-fed y no explica la ruta.
\end{remark}

\section{Extensiones proyectivas y matriciales aportadas por los entrelazadores operatoriales}

Forcing organiza posibles prolongaciones; Borel clasifica conjuntos de
historias; minimax construye robustez frente a falsadores. Ninguno sustituye
la ley estructural que produce un único superviviente. Su función es
clarificar la topología y la decisión del espacio de rutas, no introducir una
selección retrospectiva.

\begin{remark}[Separación de resultados y alcance]
Los teoremas de Gödel y Cohen son \emph{teoremas clásicos aplicados}. El
poset, la comparación tipada entre filtro de profundidad y filtro genérico,
la clase Borel HMT y el juego selector--falsador son
\emph{construcciones explícitas del presente desarrollo}. Martin, von Neumann y Sion se aplican
después, con sus hipótesis clásicas. Las antiguas afirmaciones PSC/HTT sobre
CH, AD y MM no se promocionan.
\end{remark}
